\chapter{Evaluation Design}
\label{chap:evaluation}

\section{Evaluation logic}

\begin{itemize}
  \item Evaluate whether the complete workflow turns realistic inputs into scientifically interpretable outputs.
  \item Separate four questions:
  \begin{itemize}
    \item scientific plausibility;
    \item predictive performance on hidden data;
    \item numerical/software correctness;
    \item computational feasibility.
  \end{itemize}
  \item Declare the evaluation unit for every result: fix, hidden interval, trajectory, animal, study, or computational configuration.
  \item Freeze training/test splits, configurations, and seeds before selecting final figures.
\end{itemize}

\section{Study systems and realistic interfaces}

\begin{itemize}
  \item Select case studies spanning terrestrial, marine, and airborne movement where the saved data permit defensible comparison.
  \item Include both detailed LOCAL environmental annotation and broad LONG\_RANGE land/water annotation.
  \item Candidate data/results already present for buffalo, cattle, Eastern Bettongs, boar, bears, oystercatchers, Siberian rubythroats, sharks, elk, feral cats, and whales.
  \item For each retained study, report taxon, animals, fixes, interval distribution, spatial extent, projection mode, terrain product, missingness, and pipeline role.
  \item Do not combine exploratory outputs from different code revisions until provenance and parameter equivalence have been checked.
\end{itemize}

\section{Hidden-point interpolation}

\begin{itemize}
  \item Treat high-frequency observations as partial ground truth.
  \item Hide interior points while retaining endpoints for conditioning.
  \item Evaluate regular thinning, random missing fixes, and contiguous gaps of increasing duration.
  \item Prevent leakage: hidden points must not fit states, kernels, terrain weights, scaling parameters, or tuning decisions for their evaluated interval.
  \item Metrics:
  \begin{itemize}
    \item log probability or another proper score at hidden locations;
    \item nominal contour coverage;
    \item distance to high-probability regions;
    \item invalid/no-path interval frequency;
    \item terrain/barrier violations.
  \end{itemize}
  \item Reuse and audit \code{RW-Python-gitlab/benchmark/hidden\_point\_validation/} and the mirrored results under \code{/home/omar/Desktop/MSc Thesis/Experiments/Hidden Point Validation/}.
\end{itemize}

\section{Workflow ablations}

\begin{itemize}
  \item No environmental context versus hard barriers versus relaxed/learned terrain weights.
  \item Pooled versus state-specific movement kernels.
  \item Brownian versus correlated motion.
  \item Criteria versus HMM versus BCPA versus hybrid behavioral annotation.
  \item Fixed versus automatic probability-mass kernel support.
  \item Exact UDs versus sampled-walk estimates.
  \item Native sampling interval versus each permitted Brownian time transformation.
\end{itemize}

\section{Spatial-summary validation}

\begin{itemize}
  \item Compare random-walk UDs with hidden empirical locations and aligned \code{adehabitatHR::BRB}/\code{ctmm::akde} references.
  \item Metrics: log likelihood, outside-grid percentage, zero-density percentage, integrated absolute error, Bhattacharyya overlap, contour area, and coverage.
  \item Compare candidates on the same animal and spatial support; aggregate paired differences across animals rather than across all fixes.
  \item Treat current CSV summaries as exploratory until the final manifest identifies the code revision and configuration.
\end{itemize}

\section{Terrain-weight evaluation}

\begin{itemize}
  \item Verify used and available crossing matrices independently on synthetic neighborhoods.
  \item Test whether learned weights are stable across individuals, states, raster resolutions, smoothing values, and clipping bounds.
  \item Compare learned directionality $W_{a,b}$ versus $W_{b,a}$ with held-out transitions.
  \item Evaluate whether installing the weights improves held-out prediction or reduces implausible terrain transitions relative to neutral weights.
  \item Report rare and zero-count terrain pairs instead of interpreting clipped extremes as strong evidence.
\end{itemize}

\section{Runtime, memory, and robustness}

\begin{itemize}
  \item Measure environmental acquisition, projection, segmentation, kernel fitting, terrain preprocessing, forward density, sampling, UD computation, export, and visualization separately.
  \item Vary animals, fixes, grid size, $T$, kernel support, direction bins, terrain diversity, and requested samples.
  \item Compare serial CPU, OpenMP atomic, OpenMP thread-local, and CUDA where supported.
  \item Report hardware, compiler/build flags, warm-up, repetitions, median/IQR, peak memory, transfer time, and numerical discrepancy.
  \item Stress-test invalid CRS, missing terrain, insufficient state samples, unreachable endpoints, long gaps, raster boundaries, and unavailable remote environmental services.
\end{itemize}
